\section{The actual current reduced ideal and residual order}

\begin{lemma}[Order reverses ideal inclusion]
\label{hir:order-antitone-ideal}
\uses{hir:order}
\lean{Hironaka.HironakaCurrentStrictTransformIdeal.order_le_of_ideal_le}
\leanok
If $\mathfrak a\subset\mathfrak b$ are ideals of germs, then
$\operatorname{ord}_x(\mathfrak b|_N)\leq\operatorname{ord}_x(\mathfrak a|_N)$.
\end{lemma}
\begin{proof}
If every member of $\mathfrak b_x$ belongs to the $q$th power of the
maximal ideal on $N$, the same holds for every member of $\mathfrak a_x$.
Each integer occurring in the supremum defining the order of $\mathfrak b$
therefore occurs in the supremum for $\mathfrak a$. Taking suprema proves
the assertion, including infinite orders.
\end{proof}

\begin{lemma}[Current strict transform lies in the actual cosupport]
\label{hir:current-strict-transform-in-cosupport}
\uses{hir:admissible-calculus,hir:embedded-resolution-boundary}
\lean{Hironaka.HironakaCurrentStrictTransformIdeal.strictTransform_subset_cosupp_of_centresAvoid}
\leanok
Let $W$ be closed and $W\subset\overline{W\cap O}$ for an open $O$.
Let the marked ideal $(X,X,E_0,\mathcal I_W,1)$ be in the invariant domain.
For every actual test tower $T$ whose centres avoid $T^{-1}O$, its current
strict transform $S$ is contained in its actual current cosupport.
No smoothness of $W$ is required.
\end{lemma}
\begin{proof}
The actual current ambient carrier is the whole manifold and its mark is one,
so its cosupport is the zero locus of its coherent controlled ideal and is
closed. Density transports to $S\subset\overline{S\cap T^{-1}O}$.
Every point of $S$ maps into $W$, since $W$ is closed. On $T^{-1}O$ the
actual cosupport is exactly the inverse image of $W$: restriction of the
recorded tower to $O$ has only empty centres and is a local biholomorphic
pullback. Thus $S\cap T^{-1}O$ is in the current cosupport. Closure gives
the result.
\end{proof}

\begin{lemma}[The controlled ideal vanishes on the current strict transform]
\label{hir:current-controlled-ideal-in-reduced-ideal}
\uses{hir:current-strict-transform-in-cosupport}
\lean{Hironaka.HironakaCurrentStrictTransformIdeal.controlledIdeal_le_currentStrictIdeal}
\leanok
Under the preceding hypotheses, the actual current controlled ideal satisfies
$\mathfrak a_T\subset\mathcal I_S$.
\end{lemma}
\begin{proof}
The previous lemma puts $S$ in the zero locus of $\mathfrak a_T$.
For a germ $g\in(\mathfrak a_T)_x$, coherence gives a local presentation
by holomorphic generators and expresses $g$ as a holomorphic linear
combination near $x$. The same expression makes $g$ a member of the
ideal at every sufficiently nearby point. It consequently vanishes on
$S$ near $x$, which is precisely membership in $(\mathcal I_S)_x$.
\end{proof}

\begin{lemma}[Residual ideal outside the divisor]
\label{hir:residual-ideal-off-boundary}
\uses{hir:monomial-part-along}
\lean{Hironaka.HironakaCurrentStrictTransformIdeal.nmPart_mem_iff_of_notMem_boundary}
\leanok
For a coherent ideal on the whole ambient manifold, its nonmonomial part
has exactly the same germs as the ideal at a point outside the divisor.
\end{lemma}
\begin{proof}
At such a point every condition indexed by a divisor member through that
point is empty, so the constant germ $1$ belongs to the monomial part.
A residual germ multiplied by $1$ therefore belongs to the original ideal;
restriction to the whole manifold changes no coherent ideal. Conversely
every original germ belongs to its colon by the monomial part, by ideal
closure under holomorphic multiplication. This proves both inclusions.
\end{proof}

\begin{lemma}[Removing the monomial factor preserves vanishing on a dense set]
\label{hir:residual-ideal-in-dense-reduced-ideal}
\uses{hir:residual-ideal-off-boundary,hir:monomial-part-along}
\lean{Hironaka.HironakaCurrentStrictTransformIdeal.nmPart_le_vanishingIdeal_of_dense_zeroLocus}
\leanok
Let a well-formed nonzero marked ideal have the whole manifold as its
ambient carrier, and let $n>0$. If $S$ is in its ideal's zero locus and
$S\subset\overline{S\setminus|E|}$, then its actual nonmonomial ideal is
contained in $\mathcal I_S$.
\end{lemma}
\begin{proof}
The residual ideal is coherent, by the actual SNC monomial-factorization
theorem and nonzero germs. Its zero locus is therefore closed. At a point
of $S\setminus|E|$, the preceding lemma identifies its residual germs
with original germs, all of which vanish at that point. Hence the residual
zero locus contains $S\setminus|E|$ and by density contains $S$.
The same coherent-germ propagation argument as above gives inclusion in
$\mathcal I_S$.
\end{proof}

\begin{theorem}[Actual residual ideal inside the current reduced ideal]
\label{hir:current-residual-ideal-in-reduced-ideal}
\uses{hir:current-controlled-ideal-in-reduced-ideal,hir:residual-ideal-in-dense-reduced-ideal,hir:admissible-calculus}
\lean{Hironaka.HironakaCurrentStrictTransformIdeal.nonmonomialIdeal_le_currentStrictIdeal}
\leanok
Under the current-strict-transform hypotheses, suppose additionally that
$n>0$ and every initial boundary member is disjoint from $O$.
Then the actual residual ideal of the current controlled transform satisfies
$\mathcal N(\mathfrak a_T)\subset\mathcal I_S$.
\end{theorem}
\begin{proof}
Every actual centre lies over $X\setminus O$, because the recorded tower
avoids $O$. The actual admissible-sequence producer and the boundary-support
induction put every current boundary member over either the initial
boundary or $X\setminus O$. Both miss $O$. Density transports to
$S\subset\overline{S\cap T^{-1}O}$, and this dense part misses the whole
current boundary. Thus $S\subset\overline{S\setminus|E_T|}$.
The invariant domain is preserved by the actual test tower; in particular
the current ideal is coherent and nonzero and its current divisor is SNC.
Its zero locus contains $S$ by the preceding current-ideal lemma. Apply
the residual-ideal density lemma.
\end{proof}

\begin{theorem}[Fresh reduced order is bounded by the actual residual order]
\label{hir:fresh-order-bounded-current-residual-order}
\uses{hir:current-residual-ideal-in-reduced-ideal,hir:order-antitone-ideal,bcalc:nonmonomial-order-local}
\lean{Hironaka.HironakaCurrentStrictTransformIdeal.fresh_order_le_nmOrd_geoTransformAlong}
\leanok
With the hypotheses of the preceding theorem, at every current point $y$,
\[
 \operatorname{ord}_y(\mathcal I_S)
 \leq \operatorname{ord}_y\mathcal N(\mathfrak a_T).
\]
The right hand side is the actual finite first coordinate used by the
canonical general-step invariant. In particular the freshly initialized
reduced order is finite.
\end{theorem}
\begin{proof}
For every integer $q$ such that all germs of $\mathcal I_S$ vanish to order
$q$, the proved inclusion makes all residual germs vanish to order $q$.
The actual domain and nonzero-germ theorem identify this condition with
$q$ being at most the finite residual order of the current marked ideal.
Take the supremum over those integers $q$.
\end{proof}

This comparison proves only the first order-coordinate bound. It does not
identify the higher coefficient-ideal coordinates after reinitialization,
prove their strict lexicographic descent, or exclude boundary-only branches
of the old controlled cosupport from every canonical maximum. The literal
current-singular-locus condition for an initially singular set remains open.
